

  \subsection{Space-Time Graph}
    
  \begin{definition}
    Consider a graph $G = (V, A, F)$, where $V$ is a set of vertices, and $A$ and $F$ are sets of undirected edges. For $v \in V$, denote by $\pre{v} = \{ u \in V \colon \{u,v\} \in A \}$, we extend the notion for subsets of vertices $X \subset V$ by $\pre{X} = \bigcup_{v \in X} \pre{v}$. Let $ T \colon V \rightarrow \mathbb{N}$, We say that $(G,T)$ is a time graph if and only if:
\begin{enumerate}
  \item $\{x,y\} \in A \Rightarrow \Time{y} = \Time{x} + 1$  
  \item For $x,y \in V $  $\{x,y\} \in F $ if and only if there exists $z \in V $ such both $\{x,z\},\{y,z\} \in A$. 
\end{enumerate} We will call to $T$ time-function, to $A$ arrows, and to $F$ forks. Observes that from the second requirement it follows that forks can connect only vertices that have the same time value (set by $T$). 
  \end{definition}
  \begin{remark}
    The Toom's contour graph defined in \Cref{def:contour} is, by definition of it's construction, a time-graph.
\end{remark}

  \begin{definition}[History, slice.]
    For a time graph $(G,T)$ and a vertex $u \in V$, let $G_{u}$ be the subgraph of $(V,A)$ induced by $\{ v \in V \colon \Time{v} \le \Time{u} \}$. We will denote by $\History{u}$ the 'History of $u$', which is the set of vertices that are connected to $u$ in $G_{u}$. A slice is an equivalence class of the relation $u \equiv v$ if and only if $\History{u} = \History{v}$. Notice that the time function has to be constant on a slice, to see this consider vertices  $x,y$ with $T(x) > T(y)$, thus $x \notin G_{y}$ and therefore $x \in \History{x}/\History{y}$ So they are belong to different classes. Thus, we can extend the time function and the history to slices and write for a slice $K$, $\Time{K}$ and $\History{K}$.
  \end{definition}

  \begin{claim} 
    \label{claim:disj}
    Consider slices $K_{1}, K_{2}$ such that $\Time{K_1} = \Time{K_2}$, then either $K_{1}=K_{2}$ or $\History{K_1}, \History{K_2}$ are disjoint.
  \end{claim}
  \begin{proof}
Assume a non-empty intersection, so there exists $z \in \History{K_{1}}\cap \History{K_{2}}$ such that for any $x \in K_{1}$ and $y \in K_{2}$, there exist arrow-paths $z \rightarrow x$ and $z \rightarrow y$. Thus, for any $w \in \History{K_{1}}$, one can concatenate a path $w \rightarrow x \rightarrow z \rightarrow y$ and conclude that $y \in \History{K_{2}}$.
  \end{proof}

  \begin{definition}[The graph of slice]
    \label{def:gslice}
    Let $K$ be a slice. The graph of $K$, denoted by $G_{K} = (V_{K}, E_{K})$, is the graph whose vertices are slices $R \subset \History{K}$ with $\Time{R} = \Time{K} - 1$. Two vertices are connected by an edge if and only if there is a fork whose ends belong to the slices associated with them. Namely, $S, R \in V_{K}$ are connected in $G_{K}$ if and only if there is $\{s, r\} = f \in F$ such that $s \in S$ and $r \in R$.
  \end{definition}


  \begin{claim}
    \label{claim:connected}
    $G_{K}$ is connected. 
  \end{claim}

  \begin{proof}
    Let $R, S \in V_{K}$ and consider $r,s \in V$ such that $r \in R, s \in S$. Recall that $R,S \subset \History{K}$ and therefore $r,s \in \History{K}$, Thus, by definition of the 'History' there exists a path, passing trough arrows only, from $r$ to $s$ in $\History{K}$. Consider such a path that is minimal in the number of points $x$ belong to it, with $\Time{x} = \Time{K}$. By induction on the number $k$ of such points we prove that $R$ and $S$ are connected in $G_{K}$. 

    \begin{enumerate}
      \item Base. $k=0$, In this case, the path from $r$ to $s$ lies within $G_{r} (= G_{s})$. Hence $\History{r} = \History{s}$, So $R = S$. 
      \item Induction step. There exists $y$ on the path from $r$ to $s$ such that $\Time{y} = \Time{K}$. Let $x$ be the point which precedes $y$ on the path, and $z$ the point that follows $y$. Since the values of $T$ at the two ends of an arrow differ by $1$, and since $y$ gets the maximal $T$ value in $\History{K}$ it follows that: 
        \begin{enumerate}
          \item $\{x,z\} \subset \pre{y} \Rightarrow \{x,z\} \in F$. 
          \item The slices of $x$ and $z$, say $K_{x}$ and $K_{z}$ has time value $\Time{K_{x}} = \Time{K_{z}} = \Time{K} -1$, Namely they are both members of $V_{K}$. 
        \end{enumerate}
        Since the path connecting $r,x$ in $\History{K}$ has $k-1$ vertices with $\Time{\cdot} = \Time{K}$, then by the induction hypothesis, $R$ and $K_{x}$ are connected in $G_{K}$, and by similar argument so are $K_{z}$ and $S$. On the other hand, $ \{x,z\} \Rightarrow \{K_{x},K_{z}\} \in E_{K}$. Thus, $R$ and $S$ are connected in $G_{K}$.
    \end{enumerate}
  \end{proof}


  \begin{definition}[Spanned Set]
    Consider a subset $A \subset \mathbb{R}^{d}$, and $d+1$ points $x_{1}, \ldots, x_{d+1}$. We say that $\expval*{A, x_{1}, \ldots, x_{d+1}}$ is a spanned set of $A$ if for any $i \in [d+1]$ and $s \in A$, it holds that $L(s) \le L_{i}(x_{i})$ and in addition there are $v_{1},..v_{d+1} \in V$, not necessary different, that achieve the equalities $L_{i}(x_{i}) = L_{i}(x_{i})$.  We name $x_{1}, \ldots, x_{d+1}$ the poles of $A$. Furthermore, we extend the notation to a spanned slice when $A$ is a slice of some time-graph. Notice that saying $\expval*{A, x_{1}, \ldots, x_{d+1}}$ is a spanned set means that $\Size{A} = \Size{ \{x_{1}, \ldots, x_{d+1}\}}$.
  \end{definition}
  Observe that in the case where the subset $A$ is a slice. Each of these poles can be associated with a vertex of $G$ that gives the equality. If there are multiple vertices whose give equality, associate the point with one of them arbitrarily. We will use the notion of applying functions originally defined over vertices to the poles, referring to the application over the vertices associated with them. For example $\pre{x_{i}} = \pre{v}$ for the vertex $v \in A$ which satisfies $L_{i}(v) = L_{i}(x_i)$ and therefore is associated with the pole $x_{i}$. 
 Another example is referring to $x_{i}$ as a vertex, so when saying 'For $x_{i} \in V_{k}$' or 'For $\{x_{i},u\} \in F$', we refer to the associated vertex or the edge connecting it to $u$.

  \begin{claim}
    \label{claim:prevslices}
    Consider a time graph $(G,T)$, Let $K$ be a slice in $G$, and let $\hat{K} = \expval*{K, x_{1}, .. x_{d+1}}$ be a spanned slice of $K$ such that $K \neq \History{K}$. Then for some integer $k$ there exist spanned slices $\hat{K}_{0} = \expval*{K_{0}, \cdot}, ..,\hat{K}_{k}=\expval*{K_{k}, \cdot }$ and Forks $F_{1},..,F_{k}$ such that:
    \begin{enumerate}
      \item For $i \in [d+1]$ $ \prei{x_{i}} $ belongs to the poles $\hat{K}_{0},..,\hat{K}_{k}$.  
      \item Introduce an edge between $\hat{K}_{i}$ and $\hat{K}_{j}$ if and only if one of the Forks $F_{1},..,F_{k}$ consists of a pole of $\hat{K}_{i}$ and a pole of $\hat{K}_{j}$. Then the graph formed from $\hat{K}_{0},..,\hat{K}_{k}$ is connected.  
      \item $\sum_{i=0}^{k}\Size{\hat{K}_{i}} +  \sum^{k}_{i=1}\Size{F_{i}} \ge \sum_{i=1}^{d+1}L_{i}\left( \prei{v_{i}} \right) $
    \end{enumerate}
  \end{claim}
  \begin{proof}
    First Notice that $K$ doesn't contain leaves in $(V,A)$, since a leaf is connected only to itself via arrows, and therefore its 'slice' equals to its 'History'. Thus $\prei{x_{i}}$ is not empty for $i=1,..,d+1$. Next consider the graph $G_{K}= (V_{K},E_{K})$ from \cref{def:gslice}. Since $\{ x_{i}, \prei{x_{i}} \}$ is an arrow $\prei{x_{i}}$ belongs to a slice in $V_{K}$. for $i=1,..,d+1$ denote by $R_{i}$ the slice in $V_{K}$ that contains $\prei{x_i}$. 

    By \cref{claim:connected} $G_{K}$ is connected and therefore there exist a minimal collection $\{K_{0},..,K_{k}\}$ of slices form $V_{K}$ which contains $R_{1},R_{2},..R_{d+1}$, and which is connected in $G_{K}$. We pick a set of forks $\mathbf{F} = \{F_{1},..,F_{k} \}$ which realizes a spanning tree for this collection of slices. Later we will identify edges from this spanning tree with the forks from $\mathbf{F}$.  

    For $i=1,..,d+1$ we set $\prei{v_{i}}$ to be the $i$th pole of $R_{i}$, This assures (1). The set of remaining poles for $\hat{K}_0, ..,\hat{K}_{k}$ will be equal to $\cup^{k}_{i=k}F_{i}$. This will assure (ii), We will also select poles for $F_{1}, ..,F_{k}$ to form 'spanned forks' $\hat{F}_{1}, \cdots,\hat{F}_{k}$ in such way that:  
    \begin{equation*}
      \begin{split}
        \sum^{k}_{i=0}\Size{\hat{K}_{i}} + \sum^{k}_{i=1}\Size{\hat{F}_{i}} \ge \sum_{i=1}^{d+1}L_{i}\left( \prei{v_{i}} \right) 
      \end{split}
    \end{equation*}
    This will assure (3). 

    %Now we describe the selection of poles. Since we have to span..
    Let us denote by $J_{0},..,J_{2k}$ the enumerated union $K_{0},..K_{k},F_{1},..,F_{k}$. Observes that for any choice of $(d+1)(2k+1)$ elements $z_{0}^{1},z_{0}^{2}, z_{0}^{3},..,z_{0}^{d+1}..,z_{2k}^{d+1}$ in $\mathbb{R}^{d}$, such that $z_{j}^{i} \in J_{j}$ it holds that:
    \begin{equation}
      \label{equ:spanlow}
      \begin{split}
         \sum^{k}_{i=0}\Size{\hat{K}_{i}} + \sum^{k}_{i=1}\Size{\hat{F}_{i}}  \ge \sum_{j,i}L_{i}(z_{j}^{i})
      \end{split}
    \end{equation}

Later, we are going to take the points $\{z_{j}^{i}\}$ from the intersection of $J_{j}$'s. For convenience, let us assume that for the index subset $X \subset [2k]\cup\{0\}$, the intersection $\bigcap_{j \in X}J_{j}$ returns a single point. If there is more than one, then decide on the point in a way that all the chosen points simultaneously satisfy that for any index subsets $Y, X$, such that $X \cap Y$ and they both introduce a non-trivial intersection $\bigcap_{j \in Z} J_{j} : Z\in \{X,Y\}$, then:
    \begin{equation*}
      \begin{split}
        \bigcap_{j \in X} J_{j} = \bigcap_{j \in Y } J_{j} 
      \end{split}
    \end{equation*}

    Now consider a non-trivial maximal intersection $\bigcap_{j \in X} J_{j}$, the intersection is maximal in the sense that for any $j^{\prime} \notin X$ it holds that .$\bigcap_{j \in X\cup \{j^{\prime} \}} J_{j} = \emptyset$. We claim that for any $i$ and for any such $X$, there is $t \in X$ such that $J_{t}$ is the successor of $J_{j}$ on the path to $R_{i}$ for any $j \in X/\{t\}$. Assume not. Let $x, y \in X$, and denote the paths from $J_{x}, J_{y}$ to the $i$th pole by $\expval*{J_{x}, J_{x_2}, \ldots, R_{i}}$, $\expval*{J_{y}, J_{y_2}, \ldots, R_{i}}$, such that $J_{y} \neq J_{x_{2}}$. If $y = x_{2}$, $y_{2} = x$, or $x_{2} = y_{2} \in X$, then pick other $x, y$ (if there is no such, we get an immediate contradiction). In that case, $\expval*{J_{x}, J_{y}, J_{y_{2}}, \ldots, R_{i}}$ is a path from $J_{x}$ to $R_{i}$ which is different from $\expval*{J_{x}, J_{x_{2}}, \ldots, R_{i}}$. Therefore, there are at least two paths from $J_{x}$ to the $i$th pole, namely, there is a cycle in contradiction to the minimality of $K_{0}, \ldots, K_{k}$.


We will choose the $z_{j}^{i}$ by the following process. Pick a $z_{j}^{i}$ that has not been set yet. If $J_{j} = R_{i}$, then set $z_{j}^{i} \leftarrow \prei{x_{i}}$ - we call this a type-1 assignment. Otherwise, consider the path from $J_{j}$ to $R_{i}$ (by connectivity, it has to be such a one, and by minimality, there is exactly one). Let $J_{t}$ be the successor on the path, then pick $z_{j}^{i} \in J_{j} \cap J_{t}$ - similarly, we call this assignment a type-2 assignment. Since any type-2 assignment is associated with a non-trivial intersection, and any non-trivial intersection induces a type-2 assignment for all $i$'s (since for all $i$'s, one of the $J_{j}$ on its support is a successor of the rest on their path to the $i$th pole), we have that:
    \begin{equation*}
      \begin{split}
        \sum_{j,i}L_{i}(z_{j}^{i}) &= \sum_{j,i, z_{j}^{i} \in \text{type-1}}L_{i}(z_{j}^{i}) + \sum_{j,i, z_{j}^{i} \in \text{type-2}}L_{i}(z_{j}^{i}) \\
        &= \sum_{i}^{d+1}L_{i}\left( \prei{x_{i}} \right) + \sum_{X \colon \bigcap_{j\in X}J_{j} \neq \emptyset }\left(|X|-1\right) \sum_{i}^{d+1} L_{i}\left(\bigcap_{j\in X}J_{j} \right) \\ 
        &=\sum_{i}^{d+1}L_{i}\left( \prei{x_{i}} \right) 
      \end{split}
    \end{equation*}
    Plugging it into \cref{equ:spanlow} gives the desired.
  \end{proof}

  \begin{definition}
  Let $\hat{K}$ be a spanned slice, and let $\hat{K}_{0}, \ldots, \hat{K}_{k}$ be spanned slices in $H_{K}$, and let $F_{1}, \ldots, F_{k}$ be spanned forks, such that $\hat{K}_{0}, \ldots, \hat{K}_{k}, F_{1}, \ldots, F_{k}$ satisfy the conditions in \Cref{claim:prevslices}. We call $\hat{K}_{0}, \ldots, \hat{K}_{k}, F_{1}, \ldots, F_{k}$ the \textit{predecessor decomposition} of $\hat{K}$.  
  Furthermore, we say that $\hat{K}$ is a result of a \textit{sweep-cycle}, \textit{faults}, or a \textit{logical operation} if the layer at time $T(\hat{K}) - 1$ consists of sweep-cycle gates, fault gates (whose source is $\mathcal{E}$), or logical gates.
  \end{definition}

  \begin{corollary}
    Consider a spanned slice $\hat{K}$, and let $\hat{K}_{0}, \ldots, \hat{K}_{k}, F_{1}, \ldots, F_{k}$ be the \textit{predecessor decomposition} of $\hat{K}$. 
    \begin{enumerate}
      \item If $\hat{K}$ is a result of sweep-cycle, then: 
        \begin{equation*}
          \begin{split}
            \sum_{i=0}^{k}\Size{\hat{K}_{i}} +  \sum^{k}_{i=1}\Size{F_{i}} \ge \Size{\hat{K}} + 1 
          \end{split}
        \end{equation*}
          \item If $\hat{K}$ is a result of either faults or a logical operation, then: 
        \begin{equation*}
          \begin{split}
            \sum_{i=0}^{k}\Size{\hat{K}_{i}} +  \sum^{k}_{i=1}\Size{F_{i}} \ge \Size{\hat{K}}  
          \end{split}
        \end{equation*}
      \end{enumerate}
  \end{corollary}
  \begin{proof}
    By \Cref{claim:bounded_size}
    \begin{enumerate}
      \item For $g = \varphi$ and $g = \phi \circ \mathbf{H}$, the gates apply 
    \end{enumerate}
  \end{proof}


  \ctt{ 
    Same proof works when after the application of $\phi \circ \mathbf{H}$, then $\prei{x_{i}} = \phi(x_{i})$.  
    \begin{equation*}
      \begin{split}
        & \sum_{i=0}^{k}\Size{ \phi(\hat{K}_{i}) } +  \sum^{k}_{i=1}\Size{\phi(F_{i}) } \ge \sum_{i=1}^{d+1}L_{i}\left( \phi(\prei{v_{i}}) \right) \\
        & = \sum_{i=1}^{d+1}L_{i}\left( \phi(\phi(v_{i})) \right) = \Size{K}
      \end{split}
    \end{equation*}
After the application of $\phi \circ \mathbf{H}$ the slices in the history of $X$ slice are $Z$-slices. 
It seems that the most continent way to handle the case is when we map the errors to:  
\begin{equation*}
  \begin{split}
    X_{f} & \rightarrow X_{f}Z_{\psi(f)} \\
    Z_{f} & \rightarrow Z_{f}X_{\psi(f)} \\
  \end{split}
\end{equation*}
By that, we earn that when moving the dual system the error remains the same. For that to work it seems that we have also to say something like for $\mathcal{E}^{\prime} \subset \mathcal{E}$ if things work to $\mathcal{E}$ then they also should work for $\mathcal{E}^{\prime}$. But that is not true cause $\mathcal{E}$ might contains cancellations. That is a problem. 
  }
  \begin{definition}
Let $C$ be a toric encoded circuit, let $\mathcal{E}$ be a set of faults, and consider Toom's contour $G_{\mathcal{H}}$ defined by it. We extend the slice definition into the language of $X/Z$-type error clusters as follows: We define an \textit{$X$-slice} to be a slice that is supported only on the vertices in $V_{\mathcal{H},X}$. Similarly, a \textit{$Z$-slice} is a slice that is supported only on the vertices of $V_{\mathcal{H},Z}$. As in $X/Z$-type error clusters, we define the size function to be conditioned on the type.
    \begin{equation*}
      \begin{split}
        \Size{ X\text{-slice } K } & = \sum_{i}^{d+1} \max_{v \in K} L_{i}(v)\\ 
        \Size{ Z\text{-slice } K } & = \sum_{i}^{d+1} \max_{v \in K} L_{i}(\phi(v))\\ 
      \end{split}
    \end{equation*}
  \end{definition}
  \begin{claim}
    Let $K$ be an $X$-slice, such that the logical gate preceding it is $\phi \circ \mathbf{H}$. Observes that by \cref{def:contour}, that the slices in $G_{K}$ are $Z$-slices.   
  \end{claim}



